\documentclass{jsarticle}
\begin{document}
\title{}
\author{}
\date{}
\maketitle
  体外の細胞のナトリウム カリウムポンプは、細胞内のカリウムが
  細胞外のナトリウムとの高気圧と低気圧の関係と、電子の授受により、
  カリウムが細胞外のナトリウムより多ければ、細胞内の水分は
  一定に保たれる。ナトリウムが増えた場合、細胞内の水分は、減少する。
  この理由で、浸透圧は高まる。血圧が高くなるのは、この理由である。
  ルシャトリエの平衡法則から、生活していて、ナトリウムを一定に
  とっていた場合は、食事のバランスが少しずれても、目の酸素不足は、
  大抵は大丈夫だが、食事のバランスが崩れると、カリウムで細胞内の
  還元作用が保たれていたのが、水分が細胞外に出て行くと、
  目の酸素不足が生じる。食事のバランスは、このナトリウム カリウム
  ポンプに関係している。塩分はほどほどにとっていた方が安全である。
  イオン化エネルギーから、カリウムの方が、ナトリウムよりイオン化
  されやすい。還元作用は、カリウムの方が高い。カリウムが細胞内にあると、
  電子は細胞内にとどまる。カリウムが多く取り過ぎると、食事のバランスが
  崩れると、目の酸素不足が生じる。このときに、ナトリウムをとると、細胞外
  の水分が減少して、ルシャトリエの平衡法則から、細胞内の水分と細胞外
  の水分のバランスが平衡移動して、細胞内の水分の量が一定に保たれて、
  目の酸素不足が治る。


\end{document}
